\section{Lean formalization architecture}
\label{sec:lean}

\subsection{Target repository location}

A complete implementation should live under
\begin{center}
\texttt{Logic/PrimitiveRecursiveArithmetic/Derivations/Lean/IPRADerivations/FIBER/}
\end{center}
with a small public import
\texttt{IPRADerivations.FIBER}.  The compiler metadata and source generator may
live under \texttt{Tools/}, but every accepted theorem must end in the ordinary
\texttt{Prf PRAAx} calculus.  Foundation translations belong in a separate
bridge module so the portable library remains independent of application
closures.

\subsection{Existing declarations to reuse}

The implementation should import, rather than duplicate, the following checked
components.

\begin{longtable}{p{0.35\linewidth}p{0.56\linewidth}}
\toprule
Module/declaration & Use in FIBER\\
\midrule
\endfirsthead
\toprule
Module/declaration & Use in FIBER\\
\midrule
\endhead
\texttt{ConciseIPRA.Deduction} & \texttt{Prf.subst}, \texttt{weaken}, \texttt{trans}, \texttt{exE}, logical constructors, equality substitution.\\
\texttt{IPRADerivations.Core.Basic} & \texttt{inst}, open quantifier-free induction \texttt{ind}, arithmetic symbols and congruence.\\
\texttt{Core.Semiring}, \texttt{Core.Order}, \texttt{Core.Monus} & State arithmetic, bounds, comparison, monotonicity, cancellation.\\
\texttt{Core.Pairing} & Pairing, projections, and simultaneous recursion for sequence coding.\\
\texttt{Core.Search} & \texttt{Search.mu}, \texttt{search\_spec}, \texttt{search\_found}, and \texttt{search\_minimal}.\\
\texttt{Core.BoundedSums}, \texttt{Core.MultiSums} & Boolean folds, finite potentials, and master-verifier aggregation.\\
\texttt{Core.Integers}, \texttt{Core.Powers} & Dyadic names and the adjacent-motion theorem.\\
\texttt{Core.LinearAlgebra} & Matrix observations for finite algebras and spectral clusters.\\
\texttt{IPRABridge.OpenSequents}, \texttt{ScopedSequents} & Transport of completed open interface theorems to Foundation clients.\\
\bottomrule
\end{longtable}

\subsection{Module dependency order}

\begin{enumerate}
\item \texttt{FIBER/Kernel.lean}: structural aliases for residual discharge,
      specialization, local sealing, and context replacement.  This file
      should contain no arithmetic.
\item \texttt{FIBER/Bool.lean}: Boolean operations and reflection for atomic
      equality/order predicates.
\item \texttt{FIBER/Fold.lean}: bounded conjunction/disjunction folds and
      indexing lemmas.
\item \texttt{FIBER/BoundedFormula.lean}: external syntax for a restricted
      bounded formula language, generated checker terms, and reflection.
\item \texttt{FIBER/Sequence.lean}: one canonical bounded sequence coding,
      append, lookup, concatenation, and mixed-radix alternative.
\item \texttt{FIBER/SearchAdapter.lean}: one-success adapter on top of
      \texttt{Core.Search}.
\item \texttt{FIBER/Master.lean}: least local witnesses and the
      master-certificate theorem.
\item \texttt{FIBER/Demand.lean}: capped product orders, potential, monotone
      closure, and least prefixed point.
\item \texttt{FIBER/Strategy.lean}: fixed bounded games and strategy
      realization.
\item \texttt{FIBER/Interface.lean}: metadata structures at the Lean
      metalevel and corresponding object formula constructors.
\item \texttt{FIBER/Compile.lean}: proof-producing recursion over the finite
      source graph.
\item \texttt{FIBER/Analysis/Dyadic.lean} and
      \texttt{FiniteNet.lean}: the two principal cross-domain examples.
\item \texttt{FIBER/Algebra/Clusters.lean}: stable cluster observations and
      gap-guarded labels.
\item \texttt{FIBER/Checks/*.lean}: small closed examples, \texttt{\#print
      axioms}, and bridge round trips.
\end{enumerate}

This order deliberately completes the object kernel before a large metadata
language is introduced.

\subsection{Structural kernel signatures}

The first file can use the current API almost verbatim.

\begin{lstlisting}[language={},caption={Core structural targets}]
import IPRADerivations.Core.Basic

namespace LeanProofs.IPRA.FIBER

open LeanProofs.ConciseIPRA
open LeanProofs.IPRA
open Syn
open scoped Fm

variable {n : Nat}
variable {Gamma Delta : List (Fm PTm n)}
variable {phi : Fm PTm n}

/-- Simultaneous residual discharge. -/
theorem discharge
    (h : Gamma |- [PRAAx] phi)
    (hs : forall psi, psi in Gamma -> Delta |- [PRAAx] psi) :
    Delta |- [PRAAx] phi :=
  h.trans hs

/-- Reuse a closed theorem at arbitrary terms and hypotheses. -/
theorem specialize
    {k : Nat} {theta : Fm PTm k}
    (h : ([] : List (Fm PTm k)) |- [PRAAx] theta)
    (sigma : Fin k -> PTm n) :
    Gamma |- [PRAAx] theta.subst sigma :=
  LeanProofs.IPRA.inst h sigma

/-- Certificate sealing is the native existential-elimination rule. -/
theorem seal
    {A : Fm PTm (n + 1)} {B : Fm PTm n}
    (hex : Gamma |- [PRAAx] Fm.ex A)
    (hbody : A :: Gamma.map Fm.shift |- [PRAAx] B.shift) :
    Gamma |- [PRAAx] B :=
  Prf.exE hex hbody

end LeanProofs.IPRA.FIBER
\end{lstlisting}

The ASCII notation in the distributed Lean file uses the repository's actual
Unicode/scoped notation.  The proof terms are direct aliases, so this is the
first file to elaborate and audit.

\subsection{One-success search adapter}

\texttt{Core.Search} searches for a zero of a binary term \(f(i,x)\).  FIBER
verifiers conventionally accept when \(v(i,x)=1\).  Define
\[
       f_v(i,x)=1\mathbin{\dot-}v(i,x).
\]
Then instantiate \texttt{Search.search\_found} and prove:

\begin{lstlisting}[language={},caption={Target theorem for the search adapter}]
theorem leastAccepted_found
    {v : PTm 2} {N i x : PTm n}
    (hvBit : BitValued v)
    (hi : Gamma |- [PRAAx] lt i N)
    (hacc : Gamma |- [PRAAx] inst2 v i x == num 1) :
    Gamma |- [PRAAx]
      lt (leastAccepted v N x) N /\
      (inst2 v (leastAccepted v N x) x == num 1 /\
       le (leastAccepted v N x) i)
\end{lstlisting}

Here \texttt{BitValued} should be a proved object formula, not a host
proposition silently used to rewrite values.  The implementation needs the
existing monus and parity/Boolean lemmas to show that the complement is zero
exactly when \(v=1\).

\subsection{Bounded formula syntax}

A convenient external datatype is:

\begin{lstlisting}[language={},caption={Metalevel fixed bounded formulas}]
inductive BFm (n : Nat)
  | eq   : PTm n -> PTm n -> BFm n
  | lt   : PTm n -> PTm n -> BFm n
  | and  : BFm n -> BFm n -> BFm n
  | or   : BFm n -> BFm n -> BFm n
  | imp  : BFm n -> BFm n -> BFm n
  | bEx  : PTm n -> BFm (n+1) -> BFm n
  | bAll : PTm n -> BFm (n+1) -> BFm n
\end{lstlisting}

Define \texttt{BFm.toFm} and a compiler \texttt{BFm.bit : BFm n -> PTm n}.
The main theorem is an external recursion producing a proof:

\begin{lstlisting}[language={},caption={Reflection endpoint}]
def BFm.reflect (A : BFm n) :
  ([] : List (Fm PTm n)) |- [PRAAx]
    (A.bit == num 1 <-> A.toFm)
\end{lstlisting}

The bounded-quantifier cases should call fixed fold theorems rather than place
\texttt{A.toFm} in an induction predicate.  The generated proof is accepted
because \texttt{BFm} is a Lean metatype; there is no object-level universal
truth predicate.

\subsection{Interface metadata versus object theorems}

A Lean structure may organize the generator, but proof fields must be genuine
\texttt{Prf} values.

\begin{lstlisting}[language={},caption={Schematic interface record}]
structure Interface where
  arity      : Nat
  queryOK    : PR (arity + 1)
  certBound  : PR (arity + 1)
  verifier   : PR (arity + 2)
  observe    : PR (arity + 2)
  validFm    : Fm PTm (arity + 2)
  soundFm    : Fm PTm (arity + 2)
  sound      : ([] : List (Fm PTm (arity + 2)))
                 |- [PRAAx] soundFm
  -- plus equivalence/stability and substitution equations
\end{lstlisting}

Merely declaring a field named \texttt{sound} is not a proof of a concrete
instance.  Every instance must construct it without \texttt{sorry}, an
unproved \texttt{axiom}, or a host-level semantic cast.

\subsection{Demand closure target}

Use a fixed number of coordinates in the first implementation.  Encode a
vector by nested pairing or a bounded sequence.  Define:

\begin{lstlisting}[language={},caption={Demand closure theorem shape}]
structure DemandSystem (r : Nat) where
  cap      : Fin r -> PTm n
  step     : Fin r -> PTm (n + r)
  monotone : ObjectMonotoneProof step

theorem closure_fixed
    (S : DemandSystem r) :
    ([] : List (Fm PTm n)) |- [PRAAx]
      IsLeastPrefixedPoint S (iterate S (sumCaps S))
\end{lstlisting}

For a first checked milestone, take \(r=2\).  This exercises pairing,
coordinatewise maximum, the potential argument, and open induction without
requiring a generic variable-length array library.  Generalize only after the
fixed-arity proof is stable.

\subsection{Dyadic theorem target}

The formal endpoint should avoid rational division:

\begin{lstlisting}[language={},caption={Cleared-denominator Cauchy endpoint}]
/-- a : precision and parameter |-> signed numerator -/
theorem dyadic_scaled_cauchy
    (a : ZTerm 2)
    (hstep : AdjacentUnitProof a) :
    ([] : List (Fm PTm 3)) |- [PRAAx]
      (le (var 0) (var 1) -->
       zabs (zsub (zscalePow2 a (var 0) (var 1))
                  (a.at (var 1) (var 2)))
         < pow two (monus (var 1) (var 0)))
\end{lstlisting}

The implementation should first prove the exact geometric bound
\(2^{m-n}-1\), then derive the strict inequality.  This module provides a
small but meaningful demonstration that a framework theorem closes a reported
coherence gap in the object theory.

\subsection{Compiler implementation strategy}

There are two reasonable stages.

\paragraph{Stage A: trusted Lean generator.}
Represent a well-formed finite program as an inductive Lean datatype whose
constructors contain the local \texttt{Prf} proofs.  Define
\texttt{compile} by recursion, returning a \texttt{Prf} term.  Lean's kernel
checks the result; the generator itself need not be verified as an object
program.

\paragraph{Stage B: certificate-producing external frontend.}
A Python or Lean elaborator may parse JSON metadata, synthesize primitive
recursive terms, and emit Lean source.  Its output is still checked by Lean.
The external tool is not in the trusted base.  Keep generated formulas and
node IDs deterministic so proof artifacts can be diffed and replayed.

Do not begin with an internal self-interpreter.  The diagonal obstruction and
the size of syntax-substitution proofs make that interface both impossible in
the strongest form and unnecessary for the repository's goals.

\subsection{Required checks}

Each completed module should pass:

\begin{enumerate}
\item \texttt{lake build +IPRADerivations.FIBER.<Module>};
\item a no-placeholder scan for \texttt{sorry}, \texttt{admit}, and new
      theorem-specific axioms;
\item \texttt{\#print axioms} on every public endpoint;
\item at least one open-term specialization test;
\item at least one nonempty-context test;
\item a native-scalar transfer test via \texttt{pra\_iff\_native};
\item and, for a Foundation client, an open-sequent bridge round trip.
\end{enumerate}

Lean's standard metatheoretic axioms may appear in the report.  The current
checked repository already uses \texttt{propext},
\texttt{Classical.choice}, and \texttt{Quot.sound}.

The important audit is that no new object axiom or placeholder enters the
theorem.
